The roots of $p_1$ are $u,-u-1$ for an integer $u\ge0$, so $k=-u(u+1)$. Then $p_2=(x-1)(x+u+1)$, and $p_3=x^2+x-u-1$. For $p_3$ to have integer roots $t,-t-1$, we need $u+1=t(t+1)$ with $t\ge0$. Similarly $p_5$ has integer roots $s,-s-1$ precisely when $t+1=s(s+1)$. The even-indexed polynomials automatically factor as above. Finally $p_7=x^2+x-s-1$ must fail to have integer roots. To maximize $k$, minimize $u$, hence minimize $s$. The choice $s=0$ makes $t=-1$, invalid; $s=1$ gives $t=u=1$, but then $p_7$ still has roots $1,-2$. The next choice $s=2$ gives $t=5$, $u=29$, and $s+1=3$ is not a product of consecutive nonnegative integers, so $p_7$ has no integer roots. Thus $k=-29\cdot30$, whose largest prime divisor is $\boxed{29}$.
